% 《大学物理学学习指导》第5版，第16章
% 例16.1—16.4；原书印刷227—228／PDF235—236页。
\subsection{例题：核力、粒子作用及结合能}

\sourceexample{16.1}{227}{235}
\problemtext 填出核子之间核力的三个基本特征。
\workings
原书所给的三个特征依次为
\textbf{短程性}、\textbf{饱和性}、
\textbf{（近似）电荷无关性}。
短程性指强核力主要作用于核子尺度；
饱和性表明核子的强作用主要涉及附近有限核子。
电荷无关近似对应强相互作用的同位旋对称性。
\crosscheck 核力近似电荷无关，不意味着质子之间的
电磁库仑斥力消失；真实核力仍有自旋及张量等依赖。

\sourceexample{16.2}{227}{235}
\problemtext 根据教材交换粒子的图像，填出参与电磁相互作用的
粒子及其作用媒介。
\workings
本题所强调的是\textbf{带电粒子及具有磁矩的粒子}，
电磁相互作用通过\textbf{光子}交换描述。
\crosscheck “交换光子”在散射理论中可指虚光子，
不要求过程中一定有可被单独探测的自由实光子。

\sourceexample{16.3}{227}{235}
\problemtext 写出本章讨论的三种基本粒子相互作用。
\workings
按原书三空依次为
\textbf{强相互作用}、\textbf{电磁相互作用}、
\textbf{弱相互作用}。
\crosscheck 更广义的四种基本相互作用还包括引力，
但本题要求的是原书核与粒子物理部分列出的三类，
不能凭空把三空改为四空。

\sourceexample{16.4}{227--228}{235--236}
\problemtext 中性碳-12原子质量为$12.0000\,\mathrm u$，
氢原子质量为$1.007825\,\mathrm u$，
中子质量为$1.008665\,\mathrm u$。
求碳-12原子核结合能与每核子的平均结合能。
\workings
利用中性氢原子质量计算，电子质量的贡献
与中性碳原子中的电子数近似抵消。
取质子数$Z=6$、质量数$A=12$，
\[
 \begin{aligned}
 \Delta m&=Zm_{\mathrm H}+(A-Z)m_n-M_{\mathrm C}\\
 &=6(1.007825)+6(1.008665)-12.0000\\
 &=0.098940\,\mathrm u.
 \end{aligned}
\]
使用$1\,\mathrm u\,c^2\simeq931.5\,\mathrm{MeV}$，
\[
 \boxed{E_b=\Delta mc^2\simeq92.163\,\mathrm{MeV}},
 \qquad
 \boxed{\overline E_b=\frac{E_b}{12}\simeq7.680\,\mathrm{MeV}}.
\]
\crosscheck 原书使用给定有效位数的原子质量，
不是精密核质量表；若混用裸质子质量与中性碳质量，
电子不能相消，结果将产生系统误差。
